\begin{definition}[Iwaniec--Manfredi qualitative regularity, as a hypothesis on $p$]
Let $p\in\mathbb{R}$.  We say that \emph{$p$ satisfies Iwaniec--Manfredi regularity} when the
following holds.  Let $\Omega\subseteq\mathrm{Plane}$ and let $u$ be $p$-harmonic in
$\Omega$ with weak gradient $g$ (which in particular makes $\Omega$ open).  Let
$k\in\mathbb{N}$ and $0<\beta<1$ satisfy
$k+\beta<s(p)$, and let $x_0\in\Omega$.  Then there exist a radius $r>0$ and a constant
$C>0$ such that $B_r(x_0)\subseteq\Omega$, the function $u$ is of class $C^{k}$ on
$B_r(x_0)$, and
\[
  \bigl\lVert D^ku(x)-D^ku(y)\bigr\rVert\le C\,\lvert x-y\rvert^{\beta}
  \qquad\text{for all }x,y\in B_r(x_0),
\]
where $D^ku(x)$ denotes the $k$-th Fr\'echet derivative of $u$ at $x$, a continuous
$k$-linear form on $\mathrm{Plane}$, and $\lVert\cdot\rVert$ its operator norm.  In other
words $u\in C^{k,\beta}(B_r(x_0))$.

Two features of the quantifier order are essential.
\begin{enumerate}
\item The radius $r$ and the constant $C$ are chosen \emph{after} $u$ and $g$, so they are
permitted to depend on the solution.  This is what makes the statement \emph{qualitative}:
as the paper emphasises, ``the qualitative theorem does not itself provide uniform bounds for
the corresponding norms'', and controlling these constants by the size of the solution is
precisely the difficulty that the paper's compactness argument overcomes.  A version of this
hypothesis with $r$ and $C$ independent of $u$ would make the paper's main theorem trivial.
\item The order is expressed through a H\"older exponent below the threshold: local
$C^{k,\beta}$ regularity is asserted for every pair $(k,\beta)$ with $0<\beta<1$ and
$k+\beta<s(p)$.
\end{enumerate}

This is the regularity theorem of Iwaniec and Manfredi (\emph{Regularity of $p$-harmonic
functions on the plane}, Rev. Mat. Iberoamericana \textbf{5} (1989), 1--19), which the paper
cites rather than reproves.  Their theorem gives local $C^{m,\alpha}$ regularity with
$m\in\mathbb{N}$, $0<\alpha\le1$ and $m+\alpha=s(p)$, and the form stated here is its
standard consequence: if $k+\beta<s(p)=m+\alpha$ with $0<\beta<1$, then either $k<m$, in
which case local $C^{m,\alpha}$ regularity gives local $C^{k,1}\subseteq C^{k,\beta}$
regularity, or $k=m$ and then $\beta<\alpha$, so local $C^{m,\alpha}$ regularity gives local
$C^{k,\beta}$ regularity directly; the case $k>m$ cannot occur, since then
$k+\beta>k\ge m+1>m+\alpha=s(p)$.  Thus the hypothesis stated here is strictly weaker than
the cited theorem.
\end{definition}
